% =============================================================
% Capitolo: Il Model Context Protocol (MCP)
% Progetto: MCP_ECOM — Agente e-commerce basato su MCP
% =============================================================

\chapter{Il Model Context Protocol (MCP)}
\label{chap:mcp}

% ─────────────────────────────────────────────────────────────
\section{Introduzione}
\label{sec:mcp-intro}
% ─────────────────────────────────────────────────────────────

Il Model Context Protocol (MCP) è uno standard aperto proposto da
Anthropic nel novembre 2024 con l'obiettivo di risolvere un problema
fondamentale nell'ingegneria dei sistemi basati su Large Language Model
(LLM): la \emph{frammentazione} tra modelli e sorgenti di dati.

Prima dell'avvento di MCP, ogni integrazione tra un LLM e uno
strumento esterno — che fosse un database, un'API REST o un sistema di
file — richiedeva un connettore \emph{ad~hoc}, scritto e mantenuto
separatamente per ciascun modello e ciascun provider. Il risultato era
un ecosistema a silos, costoso da mantenere e difficile da evolvere.

MCP definisce un'interfaccia universale basata sul modello
\textbf{Host / Client / Server}: il modello non comunica direttamente
con le sorgenti di dati, bensì attraverso un protocollo strutturato che
consente di esporre \emph{strumenti}, \emph{risorse} e \emph{prompt}
in modo indipendente dall'implementazione sottostante.

\begin{figure}[htbp]
    \centering
    % Placeholder: inserire diagramma "Host–Client–Server MCP"
    \fbox{\parbox{0.7\linewidth}{\centering\vspace{2cm}
        \textit{[Diagramma: Host — MCP Client — MCP Server]}\\
        \vspace{2cm}}}
    \caption{Architettura generale del Model Context Protocol.}
    \label{fig:mcp-architecture}
\end{figure}


% ─────────────────────────────────────────────────────────────
\section{Architettura del Protocollo}
\label{sec:mcp-arch}
% ─────────────────────────────────────────────────────────────

\subsection{Ruoli: Host, Client e Server}
\label{subsec:mcp-roles}

Il protocollo distingue tre ruoli distinti:

\begin{description}
    \item[Host] È l'applicazione con cui l'utente interagisce
        direttamente (ad esempio una chat, un IDE o — nel nostro caso —
        un agente ReAct). L'Host avvia le sessioni MCP e coordina uno
        o più Client.

    \item[MCP Client] Componente software che gestisce la connessione
        con un singolo Server MCP. Si occupa della negoziazione del
        protocollo (\texttt{initialize}/\texttt{initialized}), del
        catalogo degli strumenti disponibili e della chiamata ai tool.

    \item[MCP Server] Processo o microservizio che espone
        \emph{capabilities} al client: \textbf{Tools}, \textbf{Resources}
        e \textbf{Prompts}. Ogni server dichiara staticamente le proprie
        capabilities al momento dell'inizializzazione.
\end{description}

\subsection{Primitive del Protocollo}
\label{subsec:mcp-primitives}

MCP definisce tre primitive fondamentali:

\begin{enumerate}
    \item \textbf{Tools} — Funzioni invocabili dall'LLM con argomenti
        strutturati (schema JSON). Il server esegue la logica e
        restituisce un risultato testuale o strutturato. Sono la
        primitiva più usata per l'integrazione con API esterne.

    \item \textbf{Resources} — Endpoint di lettura identificati da un
        URI (anche templatizzato, es.\ \texttt{profile://\{session\_id\}}).
        Permettono all'LLM di leggere dati contestuali senza invocazioni
        esplicite di tool, analogamente a file o endpoint REST in sola
        lettura.

    \item \textbf{Prompts} — Template di sistema o di turno che il
        server può iniettare nella conversazione. Consentono di
        centralizzare le istruzioni di sistema nel server, mantenendole
        disaccoppiate dal codice dell'Host.
\end{enumerate}

\subsection{Meccanismi di Trasporto}
\label{subsec:mcp-transport}

La specifica MCP prevede due meccanismi di trasporto:

\begin{description}
    \item[\texttt{stdio}] Comunicazione tramite standard input/output:
        il server è un sottoprocesso avviato dall'Host. Adatto per
        tool locali (CLI, script, filesystem).

    \item[HTTP + Server-Sent Events (SSE) / Streamable HTTP]
        Il server è un processo HTTP separato. Il client invia richieste
        POST e può ricevere risposte in streaming tramite SSE. Questo
        trasporto è preferito per deployment distribuiti e per server
        condivisi tra più sessioni. Lo standard più recente
        (\emph{streamable HTTP}) unifica il canale bidirezionale su
        un singolo endpoint HTTP.
\end{description}


% ─────────────────────────────────────────────────────────────
\section{Implementazione nel Progetto MCP\_ECOM}
\label{sec:mcp-impl}
% ─────────────────────────────────────────────────────────────

Il progetto MCP\_ECOM adotta MCP come \emph{layer di disaccoppiamento}
tra l'agente ReAct (l'Host) e tutti i servizi e-commerce sottostanti.
La libreria utilizzata è \textbf{FastMCP} (pacchetto \texttt{mcp}),
che consente di dichiarare tool, resource e prompt con decoratori Python
analoghi a quelli di FastAPI.

\subsection{Architettura Dual-World}
\label{subsec:mcp-dual-world}

Una scelta architetturale distintiva del progetto è la presenza di
\emph{due istanze MCP indipendenti}, montate sulla stessa applicazione
FastAPI tramite ASGI sub-routing:

\begin{description}
    \item[Standard World (\texttt{/standard})] Espone tutti i tool di
        business logic (ricerca prodotti, analisi venditore, wishlist,
        ecc.) che non richiedono un browser reale. I risultati sono
        deterministici e cacheable.

    \item[Playwright World (\texttt{/playwright})] Espone i tool che
        richiedono una sessione Chromium reale (browser automation).
        Al momento dell'avvio, questo mondo eredita \emph{per
        sincronizzazione esplicita} tutti i tool non-browser dal mondo
        Standard, garantendo che l'agente in modalità browser abbia
        comunque accesso ai tool di supporto.
\end{description}

\begin{figure}[htbp]
    \centering
    % Placeholder: inserire diagramma "Dual-World MCP su FastAPI"
    \fbox{\parbox{0.8\linewidth}{\centering\vspace{2.5cm}
        \textit{[Diagramma: FastAPI Router → /standard (FastMCP) + /playwright (FastMCP)]}\\
        \vspace{2.5cm}}}
    \caption{Architettura Dual-World: due istanze FastMCP montate sullo
    stesso processo FastAPI.}
    \label{fig:mcp-dual-world}
\end{figure}

Il file \texttt{app/mcp/asgi.py} è il punto di entrata ASGI: crea le
due sub-app, ne avvia i \texttt{session\_manager} tramite un
\texttt{lifespan} asincrono e provvede alla sincronizzazione
\texttt{sync\_standard\_tools()} all'avvio.

\subsection{Modalità di Connessione del Client}
\label{subsec:mcp-client-modes}

Il componente \texttt{MCPToolClient} (\texttt{app/mcp/client.py})
implementa una doppia modalità di connessione:

\begin{enumerate}
    \item \textbf{Local Direct Mode} — Se il server URL è
        \texttt{127.0.0.1} o \texttt{localhost}, il client risolve
        direttamente l'istanza FastMCP in memoria tramite Python,
        eliminando la latenza HTTP. Questa modalità è quella utilizzata
        in produzione, poiché l'agente e il server MCP condividono lo
        stesso processo.

    \item \textbf{HTTP Remote Mode} — Se il server è remoto, il
        client utilizza \texttt{streamable\_http\_client} della libreria
        MCP ufficiale, aprendo una sessione asincrona con handshake
        \texttt{initialize}/\texttt{initialized}.
\end{enumerate}

Il client espone tre operazioni chiave, ciascuna disponibile in entrambe
le modalità:

\begin{itemize}
    \item \texttt{list\_tools\_async()} — Recupera il catalogo dei tool.
    \item \texttt{call\_tool\_async(name, args)} — Invoca un tool e
        normalizza il risultato.
    \item \texttt{read\_resource\_async(uri)} — Legge una risorsa tramite URI.
    \item \texttt{get\_prompt\_async(name)} — Recupera un prompt
        registrato nel server.
\end{itemize}

\subsection{Tool Esposti}
\label{subsec:mcp-tools}

La Tabella~\ref{tab:mcp-tools-standard} elenca i tool registrati nel
mondo Standard; la Tabella~\ref{tab:mcp-tools-playwright} quelli
esclusivi del mondo Playwright.

\begin{table}[htbp]
\centering
\caption{Tool MCP — Mondo Standard}
\label{tab:mcp-tools-standard}
\begin{tabular}{lp{0.55\linewidth}}
\hline
\textbf{Nome Tool} & \textbf{Funzione} \\
\hline
\texttt{search\_products}       & Pipeline completa di ricerca prodotti: parsing LLM della query, chiamata eBay API, ranking ibrido BM25+dense, reranking LTR. \\
\texttt{compare\_products}      & Confronto multi-prodotto per attributi, prezzo e trust score. \\
\texttt{profile\_query}         & Arricchisce la query con le preferenze dell'utente profilato. \\
\texttt{get\_item\_details}     & Recupera i dettagli completi di un articolo eBay tramite item ID. \\
\texttt{get\_shipping\_costs}   & Stima i costi di spedizione per un articolo verso un paese target. \\
\texttt{get\_similar\_items}    & Suggerisce articoli simili a quello indicato. \\
\texttt{get\_ebay\_deals}       & Recupera le offerte del giorno e promozioni attive su eBay. \\
\texttt{market\_trends}         & Analizza i trend di mercato per categoria o keyword. \\
\texttt{analyze\_seller}        & Estrae feedback, trust score e sentiment ratio di un venditore. \\
\texttt{contact\_seller}        & Contatta un venditore tramite il form eBay (modalità headless). \\
\texttt{manage\_wishlist}       & Gestisce la wishlist utente con sincronizzazione eBay. \\
\texttt{update\_user\_preferences} & Aggiorna le preferenze persistite nel profilo utente. \\
\texttt{conversation}           & Gestisce il contesto conversazionale multi-turno. \\
\texttt{ebay\_scrape}           & Estrae dati da pagine eBay tramite Playwright (headless). \\
\texttt{browser\_navigate}      & Naviga a un URL nella sessione browser aperta. \\
\texttt{browser\_click}         & Clicca su un elemento identificato da selettore CSS o testo. \\
\texttt{browser\_type}          & Inserisce testo in un campo della pagina corrente. \\
\texttt{browser\_get\_view}     & Restituisce il contenuto testuale/HTML della pagina corrente. \\
\texttt{browser\_close}         & Chiude la sessione browser. \\
\hline
\end{tabular}
\end{table}

\begin{table}[htbp]
\centering
\caption{Tool MCP — Mondo Playwright (esclusivi)}
\label{tab:mcp-tools-playwright}
\begin{tabular}{lp{0.55\linewidth}}
\hline
\textbf{Nome Tool} & \textbf{Funzione} \\
\hline
\texttt{contact\_seller\_playwright} & Automazione completa dell'invio messaggi al venditore tramite browser Chromium visibile con profilo Chrome persistente. Gestisce navigazione, compilazione form e submit. \\
\hline
\end{tabular}
\end{table}

\subsection{Resource Esposte}
\label{subsec:mcp-resources}

Il server standard espone una risorsa dinamica templatizzata:

\begin{itemize}
    \item \texttt{profile://\{session\_id\}} — Restituisce il profilo
        completo dell'utente autenticato: username, email, brand
        preferiti, budget per categoria, preferenze di condizione
        (nuovo/usato) e profondità di interazione. Il profilo viene
        arricchito dinamicamente dal servizio di \emph{user profiling}
        basato su apprendimento dalle sessioni passate.
\end{itemize}

\subsection{Prompt Esposti}
\label{subsec:mcp-prompts}

È registrato un prompt di sistema centralizzato:

\begin{itemize}
    \item \texttt{shopping\_expert\_prompt} — Istruisce l'LLM a
        comportarsi come un assistente e-commerce professionale ed
        etico. Fa riferimento esplicito all'utilizzo attivo dei tool
        di ricerca e alla lettura delle Resource utente per la
        personalizzazione delle risposte.
\end{itemize}

Centralizzare il system prompt nel server MCP — anziché nell'Host —
consente di aggiornarlo indipendentemente dalla logica dell'agente e di
riutilizzarlo in contesti diversi (web app, CLI, API).

\subsection{Contesto e Dipendenze dei Tool}
\label{subsec:mcp-context}

Ogni tool riceve un oggetto \texttt{MCPToolContext} costruito
dinamicamente a partire da:

\begin{itemize}
    \item Una \textbf{sessione database} (SQLAlchemy) aperta e chiusa
        tramite context manager, per garantire la corretta gestione
        delle risorse anche in caso di eccezione.
    \item Un oggetto \textbf{User} risolto dall'identificativo di
        sessione, per la personalizzazione dei risultati.
    \item Il parametro \textbf{llm\_engine}, che specifica il backend
        LLM da utilizzare all'interno del tool (Ollama, Gemini, ecc.).
\end{itemize}

\begin{figure}[htbp]
    \centering
    % Placeholder: inserire diagramma "Request Flow"
    \fbox{\parbox{0.8\linewidth}{\centering\vspace{2.5cm}
        \textit{[Diagramma: Agent → Executor → MCPToolClient
        → FastMCP (Local) → Tool → Service → eBay API]}\\
        \vspace{2.5cm}}}
    \caption{Flusso completo di una richiesta: dall'agente ReAct
    al tool MCP fino all'API eBay.}
    \label{fig:mcp-request-flow}
\end{figure}


% ─────────────────────────────────────────────────────────────
\section{Vantaggi Tecnici dell'Adozione di MCP}
\label{sec:mcp-benefits}
% ─────────────────────────────────────────────────────────────

\subsection{Modularità e Separazione delle Responsabilità}
\label{subsec:mcp-modularity}

L'adozione di MCP introduce una separazione netta tra \emph{orchestrazione}
(l'agente) e \emph{esecuzione} (i tool). Ogni tool è un modulo Python
indipendente — con il proprio file in \texttt{app/mcp/tools/} — che può
essere sviluppato, testato e deployato autonomamente. L'agente non ha
conoscenza dell'implementazione interna dei tool: conosce solo il loro
schema di input/output, descritto in JSON Schema.

Questa architettura consente di:
\begin{itemize}
    \item Aggiungere nuovi tool senza modificare il codice dell'agente
        (basta registrare il decoratore \texttt{@mcp.tool}).
    \item Sostituire l'implementazione di un tool (es.\ migrare da
        scraping a API ufficiale) in modo trasparente all'agente.
    \item Testare i tool in isolamento, invocandoli direttamente
        tramite la Local Direct Mode del client.
\end{itemize}

\subsection{Sicurezza e Controllo degli Accessi}
\label{subsec:mcp-security}

Il layer MCP agisce come \emph{boundary di sicurezza} naturale:
l'LLM non ha accesso diretto al database, all'API eBay o al
filesystem. Ogni operazione passa attraverso un tool con schema
validato (Pydantic/Annotated), che può applicare:

\begin{itemize}
    \item Validazione degli argomenti prima dell'esecuzione.
    \item Autorizzazione basata sul profilo utente (session\_id).
    \item Rate limiting applicabile a livello di tool.
    \item Logging strutturato di ogni invocazione per audit.
\end{itemize}

\subsection{Caching e Performance}
\label{subsec:mcp-caching}

Il layer MCP si integra con il sistema di caching Redis del progetto
(definito in \texttt{app/config/cache.py}). La chiave di cache segue
il pattern \texttt{\{mcp\_mode\}:\{tool\_name\}:\{hash(input)\}},
con TTL differenziati per tier (2 min per l'executor, fino a 24 h
per la history utente). I risultati degli strumenti browser
(\texttt{contact\_seller\_playwright}) sono marcati come terminali
e non vengono cachati, evitando loop di retry su operazioni
con effetti collaterali.

\subsection{Interoperabilità e Portabilità}
\label{subsec:mcp-interop}

Poiché il server MCP parla un protocollo standard, i tool sviluppati
per questo progetto sono in linea di principio riutilizzabili da
qualsiasi client MCP-compatibile (Claude Desktop, altri agenti,
strumenti di sviluppo come IDE plugin). Viceversa, il progetto può
integrare server MCP di terze parti — dati finanziari, cataloghi
prodotto, sistemi di pagamento — senza modificare l'architettura
dell'agente.

% ─────────────────────────────────────────────────────────────
\section{Riepilogo}
\label{sec:mcp-summary}
% ─────────────────────────────────────────────────────────────

Il Model Context Protocol costituisce la spina dorsale architetturale
del progetto MCP\_ECOM. La scelta di adottare lo standard Anthropic ha
permesso di strutturare il sistema in componenti disaccoppiati e
testabili, di gestire in modo sicuro l'accesso ai dati sensibili e di
rendere l'agente indipendente dalle implementazioni concrete dei servizi
sottostanti.

L'architettura Dual-World — con un mondo Standard per la business logic
e un mondo Playwright per l'automazione browser — è una soluzione
originale che estende il protocollo per gestire tool con requisiti
d'esecuzione eterogenei, mantenendo un'interfaccia uniforme verso
l'agente.

I capitoli successivi descriveranno nel dettaglio la pipeline di
ricerca prodotti (\S\ref{chap:search}), il sistema di analisi venditore
e la pipeline di ranking LTR.
